% Study sheet: Volume IV, Chapter 17, Responsible Physical AI.
% Numbers: \val{...} from responsibility_studysheet_numbers.tex, written by responsibility_studysheet_numbers.py.
\documentclass{studysheet}
\input{responsibility_studysheet_numbers}
\sheetvolume{4}
\sheetchapter{17}{Responsible Physical AI}
\sheetsources{Numbers from the chapter's mlsysim scenarios for the warehouse mobile manipulator (illustrative inputs); questions marked assumed use new inputs.}

\begin{document}
% ===================================================================== FRONT: LEARN
\frontside
\begin{bigidea}
A log can rebuild every tick and still not say who decided the machine could be there. Each decision needs a \hl{named owner}, each human safeguard must be one a
person can perform, and each group at risk must be measured, before the machine moves: energy
delivered across the \hl{causal boundary} cannot be recalled.
\end{bigidea}

\begin{leftcol}
\sheetsec{Key ideas}
\begin{ideas}
\idea{Reconstruction is not attribution.}{No tick names who signed the envelope; the crossing rule
  and staffing ratio sit in no record at all.}
\idea{A human role needs time, information and authority.}{Failing one adds no protection
  and moves blame onto the holder (a moral crumple zone).}
\idea{An assistant serving a fleet is a queue.}{Utilization and wait set a productivity budget,
  never a safety layer.}
\idea{Permission is granted in stages.}{Spotter, exclusion zone, open floor: each binds one
  configuration; a change, an incident or a withheld record withdraws it.}
\idea{Fairness is measured in meters and joules.}{A per-group detection gap becomes lost clearance
  and impact energy.}
\idea{Count who did not choose the risk.}{Bystander exposure, records about people and labor moved
  onto staff each get an entry.}
\end{ideas}

\vspace{9pt}
\sheetsec{The approval-record gate}
\begin{sheettable}{@{}L l@{}}
\thead{Entry missing or failing} & \thead{Gate}\\ \midrule
A decision has no owner or has expired & \textsc{refuse}\\
A role cited as a countermeasure is \textsc{infeasible} or \textsc{unknown} & \textsc{refuse}\\
Stage covers another configuration or lacks reporting duties & \textsc{refuse}\\
A group the design must detect has no per-group entry & \textsc{refuse}\\
A record used past its purpose or retention & suspend that use\textsuperscript{*}\\
Labor hours (takeovers, sweeps, monitoring) & inform review\\
\end{sheettable}
\par\vspace{3pt}{\small \textsuperscript{*}Refuse if motion cannot run without that capture. An open
entry runs under a restriction and an owner until it closes.}

\vspace{9pt}
\fcard{Leaving a stage on evidence}%
  {\ExposureDuration = \ln(1/\HypothesisTestSignificanceLevel)\times\text{target interval}}%
  {Event-free hours supporting ``under one event per target interval'' at confidence
   $1-\HypothesisTestSignificanceLevel$: $\ln(1/\val{alpha})$ = \val{lnalpha}, so one per \val{starget}\,h
   needs \val{shours}\,h, for one configuration.}%
  {}{}

\end{leftcol}%
\begin{rightcol}
\fcard{One assistant, many machines}%
  {\ServingUtilization = \ArrivalRate\,\IntervalMean, \qquad
   \QueueWaitTime = \frac{\ServingUtilization\,\IntervalMean}{1-\ServingUtilization}}%
  {Fleet request rate $\ArrivalRate$, mean time per request $\IntervalMean$, utilization
   $\ServingUtilization$, mean wait $\QueueWaitTime$.}%
  {Time test: at \val{v}\,m/s a person \val{dclear}\,m ahead leaves \val{window}\,s; an out-of-loop
   takeover needs \val{takeover}\,s.}{}
\fcard{Detection gap to impact}%
  {\begin{gathered}\ReliableDetectionFrameCount = \left\lceil\frac{\ln \DetectorMissProbability}{\ln(1-\PerFrameDetectionProbability)}\right\rceil\\[2pt]
   \DetectionRangeShortfall = \SpeedAtDetection\,\CameraFramePeriod\,\FiniteChange\ReliableDetectionFrameCount\\[2pt]
   \ImpactSpeed = \sqrt{2\CredibleBrakingDeceleration\DetectionRangeShortfall},
   \qquad \KineticEnergy = \Mass\,\CredibleBrakingDeceleration\,\DetectionRangeShortfall\end{gathered}}%
  {Frames in view so all miss with probability $\le\DetectorMissProbability$, at per-frame detection
   $\PerFrameDetectionProbability$. A group needing $\FiniteChange\ReliableDetectionFrameCount$ more frames
   loses $\DetectionRangeShortfall$, so a stop
   sized for the best-detected group reaches it still moving.}%
  {Holds while $\DetectionRangeShortfall$ is under the braking distance; misses that repeat make it worse.
   Telling \val{sa} from \val{sb} at \val{spower}\% power takes about \val{sneed} people per group.}%
  {\begin{tikzpicture}[sheet]
     \pgfmathsetmacro\xp{\val{pxp}}\pgfmathsetmacro\vz{\val{pv}}\pgfmathsetmacro\kk{\val{pk}}
     \pgfmathsetmacro\xa{\val{pxb}}\pgfmathsetmacro\xl{\val{pxbl}}
     \draw[dim] (0,0) -- ({\xp+8},0); \draw[dim] (0,0) -- (0,{\vz+6});
     \draw[muted, line width=0.9pt] (0,\vz) -- (\xa,\vz)
       plot[domain=\xa:\xp, samples=80] (\x,{sqrt(max(0,\vz*\vz-\kk*(\x-\xa)))});
     \draw[accent, line width=1.1pt] (\xa,\vz) -- (\xl,\vz)
       plot[domain=\xl:\xp, samples=80] (\x,{sqrt(max(0,\vz*\vz-\kk*(\x-\xl)))});
     \draw[dim] (\xa,{\vz-3}) -- (\xa,{\vz+3}); \draw[dim] (\xl,{\vz-3}) -- (\xl,{\vz+3});
     \node[small, anchor=south] at ({(\xa+\xl)/2},{\vz+5}) {brakes \val{dd}\,mm later};
     \draw[dim, dashed] (\xp,0) -- (\xp,{\vz+14});
     \node[small, anchor=south] at (\xp,{\vz+14}) {person};
     \fill[accent] (\xp,\val{pvimp}) circle (2);
     \node[lab, anchor=west] at ({\xp+5},\val{pvimp}) {seen \val{plow}: \val{vimp}\,m/s};
     \node[lab, anchor=west] at ({\xp+5},4) {seen \val{pref}: stops};
     \node[small, anchor=south west] at (3,{\vz+2}) {\val{v}\,m/s};
     \node[note, anchor=north] at ({\xp/2},-3) {speed vs distance, to scale};
   \end{tikzpicture}}
\end{rightcol}

% ===================================================================== BACK: PRACTICE
\backside
\begin{answers}
\ans{Finish}{$\lceil\val{lnmiss}/\val{lny}\rceil = \lceil\val{qy}\rceil$ = \val{ny} frames, so
  $\FiniteChange\ReliableDetectionFrameCount$ = \val{ydn}; \val{v} $\times$ \val{tframe} $\times$ \val{ydn} = \res{\val{ydd}\,mm};
  $\sqrt{\val{twoa}\times\val{ydd}\,\text{mm}}$ = \res{\val{yv}\,m/s};
  \val{mass} $\times$ \val{abrake} $\times$ \val{ydd}\,mm = \res{\val{ye}\,J}, \val{yratio} times the
  kneeling case. A best case: correlated misses lengthen the window (detection card).}
\ans{1}{$\ArrivalRate$ = \val{aqmach} $\times$ \val{aqrate}/h = \val{aqlammin}/min; $\ServingUtilization$ =
  \val{aqlammin} $\times$ \val{aqserv} = \res{\val{aqrho}}; $\QueueWaitTime$ = \val{aqrho} $\times$ \val{aqserv}\,min $\div$ \val{aqidle}
  = \res{\val{aqwait}\,min}. Cap: \val{qcap}/\val{aqper} $\to$ \res{\val{aqmax} machines}. No: even idle, takeover
  needs \val{takeover}\,s against \val{window}\,s (queue card).}
\columnbreak
\ans{2}{\val{lnalpha} $\times$ \val{atarget} = \res{\val{ahours}\,h}; $\div$ \val{afleet}\,h/day $\to$
  \res{\val{adays} days}. The new build is a new configuration: its count starts over (stage card).}
\ans{3}{No: plan for about \val{sneed} per group at \val{spower}\% power (\val{a3pct}\% of the set), so recruit
  the group deliberately. Meanwhile, a restriction with an owner (gate).}
\ans{4}{(a) and (d) refuse. (b) suspends that use, refusing only if motion needs the capture; (c) informs review (gate).}
\end{answers}

\begin{leftcol}
\sheetsec{Worked example}
\begin{practicebox}\boxtitle{Two frames late: what reaches a kneeling worker?}
The base runs at \val{v}\,m/s; camera frames every \val{tframe}\,ms; $\CredibleBrakingDeceleration$ =
\val{abrake}\,m/s$^2$; loaded mass \val{mass}\,kg; $\DetectorMissProbability$ = $10^{\val{pmissexp}}$.
The stop is sized for standing workers, seen in \val{pref} of frames; kneeling workers are seen in \val{plow}.
\begin{steps}
\item Frames: $\lceil\val{lnmiss}/\val{lnref}\rceil$ = \val{nref}; $\lceil\val{lnmiss}/\val{lnlow}\rceil$ =
  $\lceil\val{qlow}\rceil$ = \val{nlow}, so $\FiniteChange\ReliableDetectionFrameCount$ = \res{\val{dn}}.
\item Lost: \val{v}\,m/s $\times$ \val{tframe}\,ms $\times$ \val{dn} = \res{\val{dd}\,mm}, well inside the \val{dbrake}\,m braking distance.
\item Impact: $\sqrt{2\times\val{abrake}\times\val{dd}\,\text{mm}}$ = \res{\val{vimp}\,m/s}.
\item Energy: \val{mass}\,kg $\times$ \val{abrake}\,m/s$^2$ $\times$ \val{dd}\,mm = \res{\val{eimp}\,J}.
\end{steps}
\takeaway{A stop with no margin for the slower group turns a few centimeters of late detection into contact.}
\end{practicebox}

\vspace{10pt}
\sheetsec{Your turn: finish it}
\begin{practicebox}\boxtitle{Same stop, a group seen in \val{py} of frames}
Workers whose cartons hide the torso.\par\vspace{4pt}
\textbf{Given:} $\ln(1/(1-\val{py}))$ = \val{lny}; the reference still needs \val{nref} frames.\par
\textbf{You:} find the frames, the clearance lost, the impact speed and energy. Why is this a best case?\par\vspace{34pt}
\end{practicebox}
\end{leftcol}%
\begin{rightcol}
\sheetsec{Common mistakes}
\mistake{``The logs rebuild every tick, so we know who answers for it.''}%
  {The \val{links} decisions behind one aisle command live in no tick. If each can be retrieved with
   probability \val{plink}, the whole chain resolves only \val{plink}$^{\val{links}}$ = \val{chain} of the time.}{}
\mistake{``The remote assistant is our backup when someone steps out at the rack end.''}%
  {The hazard leaves \val{window}\,s; an out-of-loop takeover needs \val{takeover}\,s, while the machine
   travels \val{travel}\,m. The role records who will be blamed; the machine must stop itself.}{}
\mistake{``The site-wide encounter rate is low, so every worker is safe.''}%
  {If \val{encshare}\% of encounters fall at crossings used by \val{usershare}\% of workers, those
   workers carry \val{conc}$\times$ the average exposure and everyone else \val{rest}$\times$.}{}

\vspace{2pt}
\sheetsec{Check yourself}
\begin{qs}
\item \assumed{} one assistant serves \val{aqmach} machines, each asking \val{aqrate} times per hour,
  \val{aqserv}\,min per request. Utilization and mean wait? Largest fleet under a \val{qcap} cap? May the
  assistant be the countermeasure at a rack end?
\item \assumed{} the open-floor stage needs fewer than one intervention per \val{atarget}\,h at
  \val{conf}\% confidence, from \val{amach} machines running \val{aday}\,h a day. How many failure-free
  hours and days? A new enforcer build ships on day \val{aship}: what happens?
\item \assumed{} \val{a3share}\% of a \val{stotal}-person test set (\val{a3n} people) use wheelchairs.
  Is that enough to tell detection of \val{sa} from \val{sb} for them? What next?
\item Which entries refuse release? (a) a takeover role cited for rack ends, verdict \textsc{unknown};
  (b) camera traces kept past their declared period; (c) sweep hours above plan; (d) kneeling workers
  relied on, no per-group entry.
\end{qs}
\end{rightcol}
\end{document}
